\section{Introduction}
\label{sec:intro}

Large language models (LLMs) with long context windows face a critical memory bottleneck: the key-value (KV) cache grows linearly with sequence length, consuming memory proportional to $2 \times n \times H \times L \times d \times 16$ bits for $n$ tokens, $H$ attention heads, $L$ layers, and head dimension $d$.
For a 8B-parameter model at 128K context, the KV cache alone requires $\sim$32~GB in FP16---often exceeding the model weights themselves.

Existing approaches to KV cache compression fall into two categories: \emph{quantization} (reducing bits per entry) and \emph{eviction} (reducing the number of entries).
Quantization methods such as KIVI~\citep{kivi2024} and KVQuant~\citep{kvquant2024} achieve 4--8$\times$ compression, while eviction methods like H\textsubscript{2}O~\citep{h2o2023} and ScissorHands~\citep{scissorhands2024} achieve 2--5$\times$.
However, no post-training method has demonstrated $>$10$\times$ runtime memory compression with controlled quality degradation.

We present a unified framework that combines three complementary techniques:
\begin{enumerate}
    \item \textbf{Asymmetric quantization} (Layer 1): Keys at 6-bit and Values at 4-bit using rotated Lloyd-Max quantization, achieving 3.2$\times$ compression at $<$1\% perplexity increase.
    \item \textbf{Attention-aware eviction} (Layer 2): StreamingLLM-style eviction that retains attention sink tokens and recent tokens while discarding middle tokens, with eviction rates guided by a noise-floor threshold derived from quantization error.
    \item \textbf{Adaptive layer selection} (Layer 3): Outlier layers (identified via calibration) bypass compression entirely, preventing catastrophic per-layer errors.
\end{enumerate}

Our contributions are:
\begin{itemize}
    \item A \textbf{unified error bound} (Theorems~\ref{thm:quant}--\ref{thm:joint}) that decomposes the output error into quantization and eviction components, explaining why Values are more sensitive than Keys and deriving optimal eviction thresholds from the quantization noise floor.
    \item An \textbf{information-theoretic analysis} showing that our perplexity-constrained compression ratio of $\sim$13$\times$ matches the Shannon rate-distortion bound at practical coding efficiency, proving the ceiling is fundamental.
    \item The identification of a \textbf{systematic gap between perplexity and task-based evaluation} for eviction-based compression, analogous to the PSNR--VMAF gap in video compression. Under task metrics, 30$\times$+ compression is achievable.
    \item A \textbf{four-tier token classification} (Critical/Standard/Degraded/Evict) shown to be the discrete approximation to the optimal reverse water-filling bit allocation.
\end{itemize}
